\subsubsection{Harmonic decomposition of the electronic return}
\label{sec:el-espectro-retorno}

The APP--TRIT--TPK trajectory first produces the record
\(\tau_e\). The harmonic representation subsequently acts on that
record and makes explicit the information evaluated by \(K_\Omega\).
In this realization, \(\pi\) and the complex unit are the
coordinates already constructed; cyclic characters provide a
language for reading the word produced.

For \(\tau\in\mathbb R^{12}\), with indices in
\(\mathbb Z/12\mathbb Z\), set
\begin{equation}
 T_k(\tau)=\sum_{j=0}^{11}\tau_j
             \exp\!\left(-\frac{2\pi i kj}{12}\right),\qquad
 \theta_k=\frac{2\pi k}{12}.
 \label{eq:el-transformada-finita}
\end{equation}
The formula linearly extends the representation of the tritic
record to the real space of words. The correlations \(C_s\),
the windows \(A_\ell\), and the reader \(P_6\) retain their
preceding definitions on this extended domain.

\begin{theorem}[Spectral representation of the return reader]
\label{thm:el-lector-espectral}
For every real word of length twelve,
\begin{align}
 C_s(\tau)&=\frac1{12}\sum_{k=0}^{11}|T_k(\tau)|^2
                         \cos(s\theta_k),
 \label{eq:el-correlacion-espectral}\\
 \Omega(\tau)&=\frac1{12}\sum_{k=0}^{11}|T_k(\tau)|^2w_k,
 \label{eq:el-omega-espectral}\\
 w_k&=-2\cos\theta_k-4\cos(2\theta_k)-6\cos(3\theta_k),\notag\\
 P_6(\tau)&=\frac1{24}\sum_{k=0}^{11}|T_k(\tau)|^2(-1)^k.
 \label{eq:el-p6-espectral}
\end{align}
\end{theorem}
\begin{proof}
Expanding the squared modulus gives
\[
 |T_k|^2=\sum_{a,b=0}^{11}\tau_a\tau_b
           \exp\!\left(-\frac{2\pi i k(a-b)}{12}\right).
\]
The orthogonality identity
\[
 \frac1{12}\sum_{k=0}^{11}
       \exp\!\left(\frac{2\pi i k n}{12}\right)
 =\begin{cases}1,&n\equiv0\pmod{12},\\0,&n\not\equiv0\pmod{12}
   \end{cases}
\]
follows by summing a geometric progression. Multiplying the
expression for \(|T_k|^2\) by \(\exp(is\theta_k)\) and summing over
\(k\) leaves only the pairs satisfying \(a-b\equiv s\).
The resulting sum is \(C_s\); taking real parts yields
\eqref{eq:el-correlacion-espectral}. The combination
\(-2C_1-4C_2-6C_3\), already derived from the windows, gives
\eqref{eq:el-omega-espectral}. Finally,
\(C_6=2P_6\) and \(\cos(6\theta_k)=(-1)^k\) give
\eqref{eq:el-p6-espectral}.
\end{proof}

\begin{corollary}[Three modes of the central record]
\label{cor:el-tres-modos}
The word \(\tau_e=(1,1,1,-1,-1,-1)^2\) has spectral
powers
\begin{equation}
 |T_2|^2=64,\qquad |T_6|^2=16,\qquad |T_{10}|^2=64,
 \label{eq:el-potencias-espectrales}
\end{equation}
and all remaining powers vanish. Its evaluation gives
\begin{equation}
 \begin{aligned}
 \Omega(\tau_e)&=\frac{64\cdot7+16\cdot4+64\cdot7}{12}=80,\\
 P_6(\tau_e)&=\frac{64+16+64}{24}=6.
 \end{aligned}
 \label{eq:el-evaluacion-espectral}
\end{equation}
\end{corollary}
\begin{proof}
With \(z_k=\exp(-i\theta_k)\), repetition of the sextuple gives
\[
 T_k=(1+(-1)^k)(1-z_k^3)(1+z_k+z_k^2).
\]
The first factor annihilates the odd indices; the second annihilates
\(k=0,4,8\). At the three remaining indices, one obtains
\[
 T_2=4-4i\sqrt3,\qquad T_6=4,\qquad T_{10}=4+4i\sqrt3.
\]
Their squared moduli are those stated. The definition of the
weights gives \(w_2=w_{10}=7\) and \(w_6=4\); all these indices are
even. The formulas of the theorem complete the evaluation.
\end{proof}

Substitution into the return character provides a second
exact expression for its dependence on the history:
\begin{equation}
 \kappa_e=
 \frac1{12}\sum_{k=0}^{11}|T_k(\tau_e)|^2
       \left(w_k+\frac{\Delta_4}{2}(-1)^k\right)-54\alpha.
 \label{eq:el-kappa-espectral}
\end{equation}
The term \(-54\alpha\) retains the chronological provenance
of \(K_\Omega\). The other two terms evaluate the harmonic
powers with the window weights and the antipodal
correlation. The complete electronic composition
\eqref{eq:el-formula-completa} thus receives an explicit reading
of its return modes. The norm \(\sqrt3/4\), the transfer
\(\varphi/\pi^2\), the counts \(22=27-5\) and \(15=5\cdot3\),
the defect \(\Delta_4\), and the quotient \(R_{\mathrm{act}}\)
retain the constructions previously proved and brought together in
Table~\ref{tab:el-factores}.

\paragraph{Order, phase origin, and conserved information.}
The alternating word \(\tau_{\rm alt}=(1,-1)^6\) contains the
same six positive and six negative signs as \(\tau_e\).
Its transform has only \(T_6=12\), so that
\(\Omega(\tau_{\rm alt})=144\cdot4/12=48\).
This witness proves the sensitivity to order of a reader left
undetermined by the sign histogram. Its role is to delimit
the reader's domain; selecting a species requires the
corresponding section and realization.

If \((S_a\tau)_j=\tau_{j+a}\), then
\(T_k(S_a\tau)=\exp(ia\theta_k)T_k(\tau)\).
The powers and the return character are therefore invariant under
cyclic translations. The complete complex coefficients
reconstruct the word by
\[
 \tau_j=\frac1{12}\sum_{k=0}^{11}T_k(\tau)\exp(ij\theta_k),
\]
whereas their squared moduli preserve the autocorrelation.
The integral chart additionally distinguishes the event histories of
\eqref{eq:el-historias-centrales}. This hierarchy specifies the
information transmitted by each reading of the same transport.

The preceding identities have algebraic proofs over
the entire declared domain. Their reproducible verification evaluates the
witnesses in \(\mathbb Z[z]/(z^4-z^2+1)\), checks the twelve
translations, and traverses the \(4096\) binary words of length
twelve to check the identity between windows and correlations.
The three modes are components of the dodecaphasic record; the
realizations of spatially bound states additionally use
the dynamical and interaction operators of their own domain.
